\documentclass[11pt]{article}
\usepackage{amsmath,amssymb}
\usepackage{booktabs}
\usepackage{graphicx}
\usepackage{hyperref}
\usepackage{authblk}
\usepackage[margin=1in]{geometry}

\title{\textbf{Supplementary Information: Unveiling Mycobacterium marinum Cellular and Metabolic Remodeling Under Simulated Microgravity}}
\author[1,*]{Richard Barker}
\author[2]{Lynn Harrison}
\author[2]{Joseph L. Clary}

\affil[1]{Department of Agricultural and Biological Engineering, Purdue University, West Lafayette, IN 47907, USA}
\affil[2]{Department of Molecular and Cellular Physiology, LSU Health Shreveport, Shreveport, LA 71103, USA}

\date{2026}

\begin{document}
\maketitle

\section{Supplementary Tables Overview}
All empirical data tables are archived in the repository under \texttt{results\_tables/} and available via Zenodo:

\begin{itemize}
    \item \textbf{Table S1} (\texttt{Table\_S1\_raw\_counts\_matrix.tsv}): Complete raw transcript count matrix ($5,510\text{ genes} \times 9\text{ samples}$) generated by \texttt{kallisto} v0.52.0 pseudoalignment of NextSeq 550 paired-end reads.
    \item \textbf{Table S2} (\texttt{Table\_S2\_normalized\_log2cpm\_expression.tsv}): Normalized $\log_2(\text{CPM}+1)$ matrix across $4,964$ expressed genes.
    \item \textbf{Table S3} (\texttt{Table\_S3\_sample\_metadata.tsv}): Experimental metadata for NASA OSDR OSD-528.
    \item \textbf{Table S4} (\texttt{Table\_S4\_differential\_expression\_3dclinostat\_vs\_1g.tsv}): 351 significant DEGs in 3D Clinostat vs Static 1g.
    \item \textbf{Table S5} (\texttt{Table\_S5\_differential\_expression\_rpm2\_vs\_1g.tsv}): 738 significant DEGs in RPM 2.0 vs Static 1g.
    \item \textbf{Table S6} (\texttt{Table\_S6\_differential\_expression\_3dclinostat\_vs\_rpm2.tsv}): 242 DEGs comparing simulation kinematics.
    \item \textbf{Table S7} (\texttt{Table\_S7\_wgcna\_module\_assignments\_goslim.tsv}): WGCNA module assignments, standardized GOSlim descriptors, and $k_{\text{within}}$ hub scores.
    \item \textbf{Table S8} (\texttt{Table\_S8\_wgcna\_module\_eigengenes.tsv}): Module Eigengene profiles across all 9 samples.
    \item \textbf{Table S9} (\texttt{Table\_S9\_wgcna\_module\_trait\_correlations.tsv}): Pearson correlation coefficients and asymptotic $p$-values.
    \item \textbf{Table S10} (\texttt{Table\_S10\_goslim\_pathway\_enrichment.tsv}): GOSlim pathway over-representation metrics ($k/n$ overlaps and Benjamini-Hochberg FDR).
    \item \textbf{Table S11} (\texttt{Table\_S11\_full\_gene\_ontology\_enrichment.tsv}): Comprehensive GO enrichment table across Biological Process and Molecular Function.
    \item \textbf{Table S12} (\texttt{Table\_S12\_tabpfn\_loocv\_predictions.tsv}): TabPFN leave-one-out cross-validation predictions and calibrated posterior probabilities.
    \item \textbf{Table S13} (\texttt{Table\_S13\_tabpfn\_feature\_importance.tsv}): TabPFN permutation feature importance rankings.
    \item \textbf{Table S14} (\texttt{Table\_S14\_cellular\_metabolic\_model\_reactions.tsv}): Reaction-level GPR mapping across 32 enzymatic reactions.
    \item \textbf{Table S15} (\texttt{Table\_S15\_metabolic\_subsystem\_perturbation\_index.tsv}): Quantitative Subsystem Perturbation Index (SPI) rankings.
    \item \textbf{Table S16} (\texttt{Table\_S16\_pan\_microbial\_osdr\_compendium\_78studies.tsv}): Compendium of 78 microbial spaceflight datasets in OSDR.
    \item \textbf{Table S17} (\texttt{Table\_S17\_cellular\_metabolic\_model\_sbml.json}): SBML-compatible computational model schema.
\end{itemize}

\section{Supplementary Figures Overview}
All supplementary figures are archived as vector PDFs and high-resolution 300 DPI PNGs in \texttt{supplementary\_figures/}:

\begin{itemize}
    \item \textbf{Figure S1}: Systems architecture, biological model, simulation hardware, and analytical framework.
    \item \textbf{Figure S2}: Transcriptomic PCA biplot ($70.1\%$ PC1 variance) and volcano plot of 351 DEGs.
    \item \textbf{Figure S3}: Gene count distribution across 5 GOSlim co-expression modules and module-trait heatmap.
    \item \textbf{Figure S4}: TabPFN tabular foundation model benchmark under LOOCV, confusion matrix, and feature importances.
    \item \textbf{Figure S5}: Enriched GOSlim functional pathways with calibrated FAIR Blue-White-Red color fill.
    \item \textbf{Figure S6}: Intramodular hub connectivity ($k_{\text{within}}$ vs $k_{\text{total}}$) and inter-module regulatory interactome.
    \item \textbf{Figure S7}: Pan-microbial spaceflight meta-analysis landscape across 78 OSDR studies and concordance matrix.
    \item \textbf{Figure S8}: Hexagonal trajectory radar mapping phenotypic concordance across 6 physiological axes.
    \item \textbf{Figure S9}: Multi-compartment cellular cross-section and Subsystem Perturbation Index (SPI) ranking.
\end{itemize}

\section{Read Quality Control \& Pseudoalignment Convergence}
Pseudoalignment against \textit{M. marinum} M strain (NC\_010612.1) converged with high pseudoalignment efficiency across all 9 biological replicates (mean $1.21\text{ million}$ reads mapped per sample). Normalized counts were converted to log-CPM values using edgeR TMM library scaling. Zero synthetic data were utilized in any calculation.

\end{document}
